% Historical source (removed by author request): references/_OLD_Übungsaufgaben/Übungsblatt 6 - Integralrechnung/Blatt6.tex:79
\begin{ExerciseBlock}
\ExerciseHeading{5}{Integration und Symmetrie}{integral_calculus_symmetry}
Untersuchen Sie, ob das folgende Integral als gewöhnliches uneigentliches Integral existiert:
\begin{equation*}
	\int_{-1}^1 \tan\left( 3x^5-x^3+7x \right) \, dx
\end{equation*}
(Hinweis: Nachdenken spart hier unnötige Rechnungen!)

\end{ExerciseBlock}
\begin{Solution}
\begin{RemarkBox}[title=Fachliche Korrektur zur folgenden Original-PDF]
Die dort angegebene Schlussfolgerung „Integral gleich null“ ist nicht gültig.
Für $g(x)=3x^5-x^3+7x$ gilt $g(-1)=-9$, $g(1)=9$ und
$g'(x)=15(x^2-1/10)^2+137/20>0$.
Somit durchläuft $g$ im Integrationsintervall unter anderem $\pi/2$;
$\tan(g(x))$ hat dort eine einfache Polstelle.
Das gewöhnliche uneigentliche Integral existiert daher nicht.
Ungerade Symmetrie allein rechtfertigt hier keinen Integralwert.
Die folgende Handschrift bleibt als unveränderte historische Quelle erhalten.
\end{RemarkBox}
\end{Solution}
\UseExerciseSolution{integral_calculus_symmetry}
